\chapter[2025 --- Brunkow et al.]{2025: Mary E. Brunkow, Fred Ramsdell, Shimon Sakaguchi}
\label{ch:2025_Brunkow}

\section*{Main Contribution}

\textit{Identified regulatory T cells (T cells) as the immune system's security guards that prevent immune cells from attacking our own body---discovering peripheral immune tolerance.}

\section{Official Citation}

for their discoveries concerning peripheral immune tolerance

\section{Background and Historical Context}

The immune system is a marvel of intricate checks and balances, enabling robust defenses against infections while, in most cases, avoiding destructive responses against the body's own tissues. How is this balance maintained? This question has puzzled immunologists for more than a century.

\subsection{The Problem of Self-Tolerance}

Every day, our immune system protects us from thousands of different microbes trying to invade our bodies. These microbes have different appearances, and many have developed similarities with human cells as a form of camouflage. The immune system must determine what it should attack and what it should defend. If this system fails, the immune system can turn against the body's own tissues, causing autoimmune diseases.

The concept of immune tolerance---the ability of the immune system to distinguish self from non-self---was first demonstrated by Peter Medawar and Frank Macfarlane Burnet in the 1940s and 1950s, for which they shared the 1960 Nobel Prize. Their work showed that the immune system could be trained to accept foreign tissues, a phenomenon called acquired immunological tolerance.

\subsection{Central vs. Peripheral Tolerance}

By the 1990s, immunologists understood that the immune system employs two main strategies to prevent self-attack:

\textbf{Central tolerance:} This occurs in the thymus, where developing T cells (thymocytes) are tested against the body's own proteins. T cells that strongly react to self-antigens are eliminated through a process called negative selection. This was well established and considered the primary mechanism of self-tolerance.

\textbf{Peripheral tolerance:} This refers to mechanisms that operate outside the thymus to control self-reactive T cells that escape central tolerance. The existence of peripheral tolerance was hypothesized but not well understood.

Many researchers believed that central tolerance was sufficient to prevent autoimmunity. However, the immune system turned out to be more complex than they believed.

\subsection{The Suppressor T Cell Controversy}

In the 1970s and 1980, a class of T cells called ``suppressor T cells'' was proposed to play a role in immune regulation. However, the suppressor T cell concept was mired in controversy. The cells were difficult to identify, their mechanisms were unclear, and many immunologists dismissed them as an artifact. By the late 1980s, the suppressor T cell hypothesis had fallen out of favor, and many researchers considered it a failed concept.

\section{The Discovery/Contribution}

The 2025 Nobel Prize laureates made discoveries that provided critical answers to how the immune system maintains balance. Through a combination of insightful observations and carefully designed experiments, they defined CD4+CD25+FOXP3+ regulatory T cells (Treg cells) and their importance in the control of self-reactive responses.

\subsection{Shimon Sakaguchi: Discovery of Regulatory T Cells}

Shimon Sakaguchi, born in 1951 in Kyoto, Japan, was swimming against the tide in 1995 when he made the first key discovery. At the time, many researchers were convinced that immune tolerance only developed through central tolerance in the thymus. Sakaguchi showed that the immune system is more complex and discovered a previously unknown class of immune cells that prevent autoimmune reactions.

In a landmark 1995 paper in the \textit{Journal of Immunology}, Sakaguchi demonstrated that a subset of CD4+ T cells expressing the IL-2 receptor alpha chain (CD25) played a critical role in preventing autoimmune diseases. When these CD4+CD25+ T cells were removed from mice, the animals developed severe autoimmune diseases affecting multiple organs, including the stomach, thyroid, and ovaries. When the cells were reinfused, the autoimmune diseases were prevented.

This was a revolutionary finding. Sakaguchi had identified a population of T cells that actively suppressed immune responses---not by eliminating self-reactive cells, but by controlling them in the periphery. These cells were later named regulatory T cells (Treg cells).

\subsection{Mary Brunkow and Fred Ramsdell: Discovery of FOXP3}

Mary E. Brunkow, born in 1960 in the United States, and Fred Ramsdell, born in 1960 in the United States, made the second key discovery in 2001. Working at Zymogenetics in Seattle, they were studying a strain of mice called ``scurfy'' that suffered from a severe autoimmune disease. These mice had a defective gene that prevented them from forming functional regulatory T cells.

Brunkow and Ramsdell identified the gene responsible for the scurfy phenotype and named it \textit{Foxp3} (forkhead box P3). They demonstrated that mutations in this gene caused the autoimmune disease by preventing the development of regulatory T cells. Their paper, published in \textit{Nature Genetics} in 2001, provided the molecular explanation for why some mice suffered from autoimmune diseases.

\subsection{Convergence: FOXP3 as the Master Regulator}

In 2003, Shimon Sakaguchi combined the discoveries of Brunkow, Ramsdell, and his own group. He demonstrated that FOXP3 was the master transcription factor that defined regulatory T cells. Sakaguchi showed that FOXP3 was specifically expressed in CD4+CD25+ T cells and that forced expression of FOXP3 in ordinary T cells converted them into regulatory T cells.

This convergence of discoveries was transformative. It provided a molecular definition of regulatory T cells: they are CD4+ T cells that express CD25 and FOXP3. The identification of FOXP3 as the master regulator of Treg cells explained how these cells develop, function, and maintain immune tolerance.

\section{Impact on Modern Medicine}

The discovery of regulatory T cells and FOXP3 has had a profound impact on immunology and medicine. Understanding how the immune system maintains balance has opened new therapeutic avenues for autoimmune diseases, cancer, and transplantation.

\subsection{Autoimmune Diseases}

Regulatory T cells are now known to be defective or insufficient in many autoimmune diseases, including type 1 diabetes, multiple sclerosis, rheumatoid arthritis, and inflammatory bowel disease. Understanding Treg cell biology has led to clinical trials testing Treg cell-based therapies for these conditions.

\subsection{Cancer Immunotherapy}

Tumors often evade the immune system by recruiting regulatory T cells to the tumor microenvironment, where they suppress anti-tumor immune responses. Understanding this mechanism has led to strategies to deplete Treg cells in tumors or block their suppressive function, enhancing the effectiveness of cancer immunotherapies including checkpoint inhibitors.

\subsection{Transplantation}

Regulatory T cells play a crucial role in preventing rejection of transplanted organs. Clinical trials are testing whether infusing donor-specific Treg cells can promote tolerance to transplanted organs, potentially eliminating the need for lifelong immunosuppressive drugs.

\section{Legacy}

The legacy of Brunkow, Ramsdell, and Sakaguchi's work extends across immunology and medicine. Their discoveries have launched a new field of research---Treg cell biology---that continues to expand rapidly.

Understanding of regulatory T cells has revealed fundamental principles of immune regulation that apply to virtually every aspect of immunology, from infection and vaccination to cancer and autoimmunity. The identification of FOXP3 as the master regulator of Treg cells has provided a molecular tool for studying immune tolerance and developing new therapies.

The therapeutic potential of Treg cell-based approaches is enormous. Clinical trials are testing Treg cell therapies for type 1 diabetes, multiple sclerosis, organ transplant rejection, and graft-versus-host disease. In cancer, strategies to modulate Treg cell function are being combined with checkpoint inhibitors and other immunotherapies to enhance anti-tumor immune responses.

Their discoveries have thus conferred the greatest benefit to humankind, providing fundamental knowledge of how the immune system is regulated and kept in check, and opening new possibilities for treating some of the most challenging diseases.

\section{Modern Developments}

Brunkow, Ramsdell, and Sakaguchi's discoveries have led to a rapidly expanding field with numerous clinical applications. Understanding of regulatory T cells has enabled development of Treg cell-based therapies for autoimmune diseases, with clinical trials underway for type 1 diabetes (using antigen-specific Treg cells), multiple sclerosis, and inflammatory bowel disease.

In cancer immunotherapy, understanding that tumors recruit Treg cells to evade the immune system has led to combination strategies: depleting Treg cells in tumors while administering checkpoint inhibitors enhances anti-tumor responses. Anti-CCR4 antibodies (mogamulizumab) that deplete Treg cells are approved for certain lymphomas and are being studied in solid tumors.

Understanding of FOXP3 has enabled development of ``super-Treg'' cells---engineered regulatory T cells with enhanced suppressive function---for treating severe autoimmune diseases and preventing transplant rejection. CAR-Treg cells, which combine chimeric antigen receptor technology with regulatory T cells, are in preclinical development for inducing antigen-specific tolerance in transplantation and autoimmune diseases.

The field continues to expand with new discoveries about Treg cell subsets, tissue-resident Treg cells, and the metabolic requirements for Treg cell function. Understanding of how Treg cells maintain stability---and how they can become unstable in disease---is informing the development of safer and more effective Treg cell-based therapies.
